\documentclass{article}

\usepackage{geometry}
% TODO: 1 inch margins as required by project spec.
\newgeometry{vmargin={15mm}, hmargin={17mm, 17mm}}
\begin{document}

\title{CITS4403 Project: Modelling the Perth Housing Crisis}
\author{Lauren Pudney}

\maketitle

\section{Background}
\label{sec:introduction}
% - Establishes the problem, its significance and a focused research question or modelling objective.
% - Summarises relevant studies and places the project in the context of existing literature.
Housing instability currently exists as a critical socioeconomic issue within 
Australian society. Despite high levels of economic prosperity, escalating housing
pressures and persistent homelessness continue to affect Australians
\cite{absEstimating, aihw2023australia}, placing disproportionate 
burden on marginalised and lower-income groups, such as
young people, migrants and renters~\cite{daley2018housing}. 
\\

Over the past two decades, housing and rental costs in Australia have increased
significantly relative to incomes, with home-ownership rates for all age
groups under 65 decreasing, particularly among the younger demographic and lower 
income households~\cite{daley2018housing}. In addition to these general affordability issues,
multiple interest rate increases by the Reserve Bank of Australia (RBA) have 
heavily increased mortgage repayments, creating more financial stress for those lucky enough
to enter the housing market~\cite{commbank2026rba}. 
\\

The purpose of this model is to act as a controlled environment to discover and analyses the
factors that drive housing instability in Perth. While the majority of studies regarding the housing market
are typically based within an Australian context, our model is specific to the Perth CBD landscape.
%- Not really happy with how the last paragraph ends

\section{Research aims}
% - Establishes the problem, its significance and a focused research question or modelling objective.
% - Summarises relevant studies and places the project in the context of existing literature.

\subsection{Research hypothesis}
Our hypothesis for this study is to determine which factors have the greatest
impact on home ownership rates and housing instability in Perth. Specifically, 
how factors such as the recent interest rate hikes, the upfront deposit amounts 
required, and the general income inequality between agents affect housing 
accessibility. Ultimately, we plan to use our model to evaluate these variables
to discover which factor has the most significant influence on decreasing Perth's 
current housing instability.
\\ % ahfbsrghb need to reword last bit, i dont like it

The reasons we chose to investigate these factors is;
\begin{itemize}
  \item Interest Rate Hikes: have been frequently in the news since the RBA has
  raised the cash rate to its highest since 2011 \cite{abcratehike}. Due to the articles about
  how home owners face higher mortgage repayments due to this increase, we wanted to test
  if the interest rate hikes actually are a major factor in Perth's housing instability
  crisis.
  \item Upfront Deposit Requirements: The government has implemented programs such as Keystart
  to reduce the required deposit amount needed to buy a house, from 20\% of its purchase
  price to 5\%. However recent news articles state that these low deposit schemes
  actually hinder homeowners, as they now have higher mortgage repayments \cite{newsdeposits}.
  We want to test if these programs are a significant factor and ultimately if 
  they improve access into the housing market.
  \item Income inequality: The wealth disparity heavily impacts the housing market,
  as the richer households can afford to buy up more housing, and then use those houses
  as rentals to generate passive income. This is turn creates a cycle of renters
  needed to rent forever due to affordable housing being brought up as investment
  properties. We want to evaluate how this wealth gap drives the overall housing 
  instability in our model, and determine if it is the most significant factor in
  the decrease of home ownership rates in Perth.
\end{itemize}


\section{Prior work and Contributions}
Many prior works attempt to simulate housing markets across varying 
contexts and approaches. Housing crises in London~\cite{kaveh2022london} and Tokyo~\cite{tamama2026housing} are 
examined via the implementation of system dynamics models and agent
based models respectively, with other works focusing on the 
design of targeted interventions for such scenarios~\cite{daniel2025data}.
Models developed in the Australian context often seek to support policy
decisions. Economic modelling approaches are used
in a 2025 paper to estimate housing need in the Australian market~\cite{daniel2025data},
whilst a 2020 model of the Australian housing market seeks to
examine relationships between construction, vacancies, rents and prices,
via an empirical modelling method,
finding low interest rates and high immigration as key factors influencing
rapid growth in housing prices and construction~\cite{saunders2020model}.
Simple agent based models simulate generic housing markets, with 
J M Applegate presenting an extendible agent based model, simulating
interactions between households buying and selling houses of varying quality~\cite{applegate2020general}.

Whilst models exist for the simulation and examination of both housing
crises, and the Australian market, we seek to combine these elements, with
a specific focus on the Perth metropolitan region. Similarly to Saunder's
2020 paper~\cite{saunders2020model}, we aim to identify relationships involving interest rates,
rates and prices, however aim to do so with specific focus on outcomes
for renters. Uniquely, we aim to simulate the Perth market via an
agent-based model, similar to that presented by Applegate, however
with significant extensions and assumptions tailored to our context.
Our approach allows us to simulate specific interactions between renters and
landlords, offering a visual representation of the distribution of
household rental status and characteristics across the city, alongside
renter specific outcomes.
\\
\\
We make the following contributions:
\begin{enumerate}
    \item We present an \textbf{agent-based model} of the housing market in the
    Perth metropolitan region
    \item Via simulation, we investigate the \textbf{impact on renter outcomes
    associated} with interest rate hikes, upfront deposit amounts 
    required, and the general income inequality within Perth.
\end{enumerate}

% - Clearly distinguishes the team’s contribution from existing approaches or taught examples.
% - Explains the purpose of that contribution and demonstrates its added value or insight


% https://onlinelibrary.wiley.com/doi/full/10.1111/1475-4932.12537

% % This is from a not-very-reputable journal and has a dodgy citation.
% Data-Driven Modelling of Migration, Income Trends, and Rental Pressures to
% Design Targeted Interventions Reducing Housing Instability~\cite{Daniel2025data}
% \\
% Modelling housing need in Australia to 2025~\cite{Rowley2017modelling}
% \\
% London's Housing Crisis; A System Dynamics Analysis of Long-term
% Developments: 40 Years into the Past and 40 Years into the Future
% \\
% HousingABM Japan: An Agent-Based Model of Price and Rent
% Amplification in the Tokyo Housing Market \\
% https://www.comses.net/codebases/0fbece58-7669-4204-870f-1643e7821fb5/releases/1.0.0/
% (not peer reviewed)
% \\


\section{Model specification}
% - Defines the components, variables, interactions, equations or update rules needed to understand the model.
% - States assumptions and initialisation rules, and justifies relevant model-design and parameter choices.
\textbf{Classes}
\\

\texttt{Dwellings}: Represents the houses in our model. Each \texttt{Dwelling}
is initialised with a tuple representing a set of coordinates (\textit{coord}),
which then determines it's \textit{quality} based on its location to hotspots, such
as the Perth CBD. It stores the current \textit{occupant} and \textit{owner}, (None
when initialised), alongside whether it is for sale or not (\textit{isOnMarket}). 
Lastly, each \texttt{Dwelling} has a rent and sale price, formulated using the 
\texttt{AVERAGE\_HOUSE\_PRICE} (or \texttt{AVERAGE\_RENT} if it's being used as 
a rental), multiplied by it's quality score, and then further multiplied by a 
random factor to account for real world variations.
\\

\texttt{Household}: Represents the people occupying the houses in our model. They're
initialised with an income using a lognormal distribution, with the mean being the
\texttt{AVERAGE\_INCOME} variable, and their cash assets (\textit{money}) are set to
double their income, so the household can start buying/renting houses immediately
(if applicable). A \texttt{Household} can occupy a house (\textit{occupying}), and
can own multiple houses and mortgages (\textit{owns} and \textit{mortgages} 
respectively).
\\

\texttt{Bank}: Represents the bank. Each \texttt{Bank} is initialised with a base
interest rate (\textit{interestRate}), and tracks all the active mortgages 
(\textit{activeMortgages}) occuring in the model. It also hosts an method to update
the interest rate each year to simulate the rate changes, with the interest rate 
unable to go above 17.5\% or below 0.1\% based on historical data \cite{ratehistory}. 
\\

\texttt{Mortgages}: Holds information about each households' mortgage 
(or mortgages). It tracks the household who is responsible for paying off the loan
(\textit{borrower}), and the property that is tied to the mortgage (\textit{dwelling}).
It also holds the \textit{initialLoanAmount} for the loan, along with the year's 
\textit{interestRate}, how many years the loan has been active for (\textit{yearsPaid}),
and the remaining balance on the loan (\textit{loanAmount}). The \texttt{Mortgages}
class also has a method to calculate the annual payments a household has to make
towards their mortgage by using the annuity formula, which most Australian banks 
implement~\cite{wikimortgage}\cite{moneysmart}.

% how helpers work
% how tick() works
% do we need to go into depth about how the graphing works?? (tbh i'm scared of that file lol)

\subsection{Assumptions}
\textbf{Income}: Median household income (weekly): \$1,865 (ABS 2021 census so slightly out of date)
\cite{abs2021greater}
\\
\textbf{Housing Prices}: Median price data: \$965,000 \cite{reiwa}.
\\
\textbf{Rental Prices}: Based on the median rental prices  (weekly): \$750 \cite{reiwa}.
\\
\textbf{Interest Rate}: Based on the cash rate implemented on the 30th of Sept 2026 
\cite{ratehistory}: 4.6\%. Will change during experimentation to determine it's effect
on home ownership using historical rates (ranging from 1\% to 17.5\%).
\\
\textbf{Profit Margin}: Assuming that home investors will sell their property once it is 
valued at 150\% its original cost, with the 150\% figure chosen arbitrarily.
\\
\textbf{Average Loan Period}: Based on the average loan period for most mortgages \cite{loanperiod}:
30 years.
\\
\textbf{House Deposit}: Based on the average home deposit in Australia \cite{deposit}: 20\%. 
Will change during experimentation to 2\% to simulate government programs such as Keystart
\cite{keystart}.
\\





\section{Methodology}
% aka experimental design
% - Uses appropriate experiments or case studies to investigate the research question and explore model behaviour.
% - Specifies scenarios, parameter ranges, outcome measures and, where relevant, repeated runs or comparisons.
% - Provides sufficient procedural detail to repeat the experiments.

\section{Results}
% - Reports relevant qualitative and quantitative findings accurately and provides sufficient evidence of model behaviour under the investigated conditions.
% - Uses well-chosen figures and tables with informative captions, labels and units where applicable.

\section{Discussion}
% - Interprets findings in relation to the research question, rather than merely repeating the results.
% - Draws evidence-supported conclusions and explains the insights gained about the system.
% - Identifies limitations, uncertainty and potential improvements or applications without overstating the model’s scope

% - limitation idea, we used a lognormal distrib when irl its more right skewed??
% - also we only update the interest rate each year, normally done quarterly??
% - typically formula is used for monthly payments - could have a different effect
% - only houses, not accounting for units which are cheaper upfront
% - when changing deposit to 5%, model doesn't include lenderes mortgage insurance, which would be an added cost
% - also doesn't account for first home buyers 
% - we don't consider other loan types, such as guarantor loans & bank of mum and dad
% - lots of factors, such as debt (credit scores), comparing home loans, building inspections, settlement & stamp duty, insurance
% - for above point \url{https://moneysmart.gov.au/home-loans/buying-a-house}

\section{Conclusion}
% same as Discussion

\section{Future work}

% \subsubsection{Credit statement}
% \emph{Pudney}: Core + solver backend implementation, benchmarking and survey design and implementation, document preparation. \emph{Grinbergs}: loss backend implementation, documentation preparation. \emph{Daggitt}: Conceptualisation, ITP backend and documentation preparation.

\bibliographystyle{splncs04}
\bibliography{bibliography}
\end{document}
